\section{The multiplication table problem}
    \label{sec:multiplication-table-problem}

Recall that the multiplication table problem asks for the order of magnitude of
\begin{equation*}
    M(N) = \# \set{ ab}{ a, b \in [N]}.
\end{equation*}
This is intimately related to the problem of estimating the proportion of integers with a divisor in a given interval. 
For any positive integer $n$, let $\tau(n, y, z)$ denote the number of divisors of $n$ in the interval $(y, z]$, and let
\begin{equation*}
    H(x, y, z) = \#\set{ n \le x}{\tau(n, y, z) \ge 1}
\end{equation*}
be the number of integers $n \le x$ that have a divisor in $(y, z]$.

The study of the special case $z = 2y$ was initiated by Besicovitch \cite{besicovitch} in the 1930s, and estimates for $H(x, y, 2y)$ were given by Erd\H{o}s \cite{anasymptoticinequality} and Tenenbaum \cite{tenenbaumsurlaprobabilite}.
In the seminal paper \cite{thedistributionofintegers}, Kevin Ford determined the order of magnitude of $H(x, y, z)$ for all values of $x, y, z$.
In this dissertation, we will focus on the case $z = 2y$, which is relevant for the multiplication table problem.

\begin{theorem}[{\cite[Theorem 1]{thedistributionofintegers}}] \label{thm:integers-with-a-divisor-in-y-2y}
    Let $x > 0$. 
    We have, uniformly for all $3 \le y \le \sqrt{x}$,
    \begin{equation*}
        H(x, y, 2y) \asymp \frac{x}{(\log y)^\delta (\log \log y)^{3/2}},
    \end{equation*}
    where $\delta = 1 - (1 + \log \log 2) / \log 2 \approx 0.086071$.
\end{theorem} 
Note that the conclusion of the theorem is trivial if $y \le y_0$ for some fixed constant $y_0$.
Indeed, in this case, if $a$ is any integer in $(y, 2y]$, then the number of multiples of $a$ in $[1, x]$ is $\asymp x$.
Henceforth, we always assume that $y$ is large.

We first show how Theorem \ref{thm:multiplication-table-problem}, which is restated here, follows from the above result.
\begingroup
\def\thedefinition{\ref*{thm:multiplication-table-problem}}
\begin{theorem}
    We have
    \begin{equation*}
        M(N)  \asymp \frac{N^2}{(\log N)^\delta (\log \log N)^{3/2}}.
    \end{equation*}
\end{theorem}
\addtocounter{definition}{-1}
\endgroup

\begin{proof}
    For the lower bound, observe that if an integer $n \le N^2 / 4$ has a divisor in $(N/4, N/2]$, it is counted by $M(N)$. 
    Hence, 
    \begin{equation*}
        M(N) \ge H\left( \frac{N^2}{4}, \frac{N}{4}, \frac{N}{2} \right) \gg \frac{N^2}{(\log N)^\delta (\log \log N)^{3/2}}.
    \end{equation*}

    For the upper bound, suppose that $n = m_1 m_2$ with $m_1, m_2 \le N$. 
    Then $m_1 \in (2^{-k-1} N, 2^{-k}N]$ for some integer $k \ge 0$, and so $n \le 2^{-k} N^2$.
    This gives
    \begin{equation*}
        M(N) \le \sum \limits_{k \ge 0} H \left(\frac{N^2}{2^k}, \frac{N}{2^{k + 1}}, \frac{N}{2^k} \right).
    \end{equation*}
    By Theorem \ref{thm:integers-with-a-divisor-in-y-2y}, the contribution from $k$ with $N/2^{k} \ge \sqrt{N}$ (say) is 
    \begin{equation*}
        \ll \left( \sum \limits_{k} \frac{N^2}{2^k} \right) \frac{1}{(\log N)^\delta (\log \log N)^{3/2}} \ll \frac{N^2}{(\log N)^\delta (\log \log N)^{3/2}},
    \end{equation*}
    since $\log (N/2^{k + 1}) \asymp \log N$ in this range.
    For the contribution from $k$ with $N/2^{k} < \sqrt{N}$, we can use the crude bound 
    \begin{equation*}
        H\left(\frac{N^2}{2^{k}}, \frac{N}{2^{k + 1}}, \frac{N}{2^{k}} \right) \le \frac{N^2}{2^{k}} \ll N^{3/2}.
    \end{equation*}
    Summing over $\frac{\log N}{2\log 2} < k \le \frac{2\log N }{\log 2}$, this gives a total contribution of $\ll N^{3/2} \log N$, which is negligible compared to the main term.
\end{proof}

When studying $H(x, y, z)$, a simple idea one might have is to write down the inequality
\begin{equation} \label{eq:trivial-upper-bound-sum}
    H(x, y, z) \le \sum \limits_{n \le x} \tau(n, y, z).
\end{equation}
By double counting, 
\begin{equation*}
    \sum \limits_{n \le x} \tau(n, y,z) = \sum \limits_{y < d \le z} \lflr \frac{x}{d} \rflr  = x \sum \limits_{y < d \le z} \frac{1}{d} + O(z) = x \log \left( \frac{z}{y} \right) + O(z).
\end{equation*}
In the range $z \asymp y$, this is useless, as it only gives an upper bound of $\ll x$.
However, when $z$ is close to $y$, very few integers have more than one divisor in $(y, z]$, and so the inequality \eqref{eq:trivial-upper-bound-sum} leads to a sharp bound.
For details, we refer the reader to \cite[Chapter 2]{divisors}.
% Intuitively speaking, the case where $z$ is close to $y$ is easier, as divisibility conditions by distinct $d, d' \in (y, z]$ are relatively independent, with $\gcd(d, d')$ being small in most cases.
% When $z \asymp y$, they exhibit more complicated dependences.

The remainder of this chapter is devoted to discussing Ford's proof of Theorem \ref{thm:integers-with-a-divisor-in-y-2y}.
The presentation mainly follows \cite{integerswithadivisor} and the lecture notes \cite{fordlecturenotes}.
The proof techniques are elementary and combinatorial in nature, although the details can be complicated. 
In the next section, we start by giving a heuristic argument for the theorem that explains some of the key ideas behind the proof.

\begin{remark}
    Similarly to the prime factorisation of integers, permutations possess a unique decomposition into cycles, and there is a well-known analogy between divisors of integers and fixed sets of permutations (cf. \cite{arratiabarbourtavare}).
    Roughly speaking, a fixed set of size $k$ corresponds to a divisor of an integer in the range $(e^k, e^{k + 1}]$.
    
    To make the connection clearer, let $C_j(\sigma)$ denote the number of cycles of length $j$ in a permutation $\sigma$.
    For integers $c_1, \dots, c_n \ge 0$ with $c_1 + 2c_2 + \dots + nc_n = n$, Cauchy's formula (see, e.g., \cite[Prop. 1.3.2]{enumerativecombinatorics}) states that the probability that a uniformly random permutation $\sigma \in S_n$ satisfies $C_j(\sigma) = c_j$ for all $j \in [n]$ is given by
    \begin{equation*}
        \prod \limits_{j = 1}^n \frac{1}{c_j!} \left( \frac{1}{j} \right)^{c_j}.
    \end{equation*}
    This suggests that the quantities $C_j(\sigma)$ behave roughly like independent Poisson variables with mean $1/j$.
    Analogously, recall from Section \ref{subsec:kubilius-model} that if $n$ is a random integer, we expect the number of prime factors in $(e^j, e^{j + 1}]$ to have an approximate Poisson distribution with mean
    \begin{equation*}
        \sum \limits_{e^j < p \le e^{j + 1}} \frac{1}{p} \approx \log (j + 1) - \log j \approx \frac{1}{j}.
    \end{equation*}
    This suggests that one might compare the probability of having a cycle of length $j$ to the probability of having a prime divisor in $(e^j, e^{j + 1}]$, or the probability of having a fixed set of size $k$ to the probability of having a divisor in $(e^k, e^{k + 1}]$.
    
    The analogous problem to Theorem \ref{thm:integers-with-a-divisor-in-y-2y} for permutations is then to determine the number of permutations in $S_n$ that have a fixed set of size $k$.
    Denote this number by $\tilde{H}(n, k)$. 
    The order of magnitude of $\tilde{H}(n,k)$ was determined by Eberhard, Ford, and Green \cite{eberhardfordgreen}, who showed that
    \begin{equation*}
        \tilde{H}(n, k) \asymp \frac{n!}{k^\delta (1 + \log k)^{3/2}},
    \end{equation*}
    uniformly for $1 \le k \le n/2$.
    In other words, $\tilde{H}(n, k)$ exhibits the same behaviour as $H(x, y, 2y)$ if one equates $y = e^{k}$.
    The proof uses the same ideas as in \cite{integerswithadivisor}, although the details are considerably simpler in the context of permutations.
\end{remark}

\subsection{A heuristic argument}
    \label{subsec:heuristic}
Let us first explain where the exponent $\delta$ comes from.
In Section \ref{subsec:multiplication-table-problem}, we already established an upper bound on $M(N)$ of the form
\begin{equation*}
    M(N) \ll \frac{N^2}{(\log N)^{Q(\lambda)}},
\end{equation*}
where $Q(\lambda) = \lambda \log \lambda - \lambda + 1$.
The optimal value of $\lambda$ turned out to be $1/\log 2$, \mbox{and $Q(1/\log 2) = \delta$.}
This gave us a hint that the main contribution to $M(N)$ comes from integers with roughly $\log \log N / \log 2$ distinct prime divisors. 
Similarly, we will see that the main contribution to $H(x, y, 2y)$ comes from integers with about $\log \log y/ \log 2$ prime divisors $\le 2y$ which, in turn, will give rise to the exponent $\delta$.
Note that by Theorem \ref{thm:hardy-ramanujan-generalisation}, a typical integer in $[1, x]$ has very close to $\log \log y$ prime divisors $\le 2y$, so similarly to the case of the multiplication table, the dominant contribution to $H(x, y, 2y)$ comes from ``atypical'' integers.

Heuristically, this can be seen as follows (cf. \cite{integerswithadivisor}).
Suppose that an integer $n \le x$ is chosen uniformly at random and write $n = ab$, where all prime factors of $a$ are $\le 2y$, and all prime factors of $b$ are $> 2y$.
Then it only depends on $a$ whether $n$ has a divisor in $(y, 2y]$.
It is reasonably likely that $y^{1/100} \le a \le y^{100}$ (say), in which case $\log a \asymp \log y$.
For simplicity, assume that $a$ is squarefree.

The correct scale to look at divisors is a logarithmic scale.
The interval $(y, 2y]$ then corresponds to an interval of logarithmic length $\log 2$.
Suppose for now that the set $\set{\log d}{d \mid a}$ is evenly distributed over $[0, \log a]$. 
Put $k = \omega(a)$. Since $\tau(a) = 2^k$, we would expect $\prob{\tau(a, y, 2y) \ge 1}$ to be roughly
\begin{equation*}
    \approx \min\left(1, \frac{2^k \log 2}{\log a}\right) \asymp \min\left(1, \frac{2^k}{\log y}\right).
\end{equation*}
There is a transition at
\begin{equation*}
    k_0 := \lflr \frac{\log \log y}{\log 2} \rflr,
\end{equation*}
so we expect the main contribution to $H(x,y,2y)$ to come from those $n$ for which $\omega(a) \ge k_0 - O(1)$.
Theorem \ref{thm:prime-factors-in-a-set} tells us that the number of integers $\le x$ with precisely $k$ prime factors $\le 2y$ is 
\begin{equation*}
    \ll \frac{x}{\log y} \frac{(\log \log y)^k}{k!}.
\end{equation*}
Summing over $k \ge k_0$, we estimate that $H(x, y, 2y)$ should be
\begin{equation*}
    \ll \frac{x}{\log y} \sum \limits_{k \ge k_0 - O(1)} \frac{(\log \log y)^k}{k!} \ll \frac{x}{\log y} \frac{(\log \log y)^{k_0}}{k_0!} \ll \frac{x}{(\log y)^\delta (\log \log y)^{1/2}},
\end{equation*}
where the last step follows from Stirling's formula.

Comparing this with Theorem \ref{thm:integers-with-a-divisor-in-y-2y}, we see that we are off by a factor of $\log \log y$.
The reason is that our uniformity assumption on the set of divisors is only justified with probability about $1/\log \log y$.
In reality, fluctuations of prime divisors on a log-log scale cause the divisors of $a$ to cluster rather than being evenly distributed, which reduces the likelihood of having a divisor in a particular interval. 

This phenomenon is related to the distribution of uniform order statistics. 
Given $X_1, \dots, X_k$ sampled independently and uniformly at random from $[0, t]$, their \defn{order statistics} are the random variables $0 \le \xi_1 \le \dots \le \xi_k \le t$, where $\xi_j$ is $j$th smallest of the sample values of $X_1, \dots, X_k$.

Suppose $a = p_1 \cdots p_k$, where $p_1 < \dots < p_k \le 2y$, and $k \sim k_0$.
In light of the Kubilius model, we expect the sequence $\log \log p_1, \log \log p_2, \dots$ to behave locally like a Poisson point process with intensity 1. (Of course, this is not true for small prime divisors if taken too literally.) Conditioned on having precisely $k$ points in $[0, \log \log (2y)]$, the law of the Poisson process is equal to the distribution of uniform order statistics on $[0, \log \log (2y)]$, see e.g. \cite[Chapter 2.3]{ross}.
Standard results on order statistics then tell us that we can expect $\log \log p_j \sim \frac{j}{k} \log \log y \sim j \log 2$, but with high probability, there exists some $j$ for which
\begin{equation*}
    \log \log p_j \le j \log 2 - c \sqrt{\log \log y},
\end{equation*}
where $c > 0$ is some small constant. (This is also roughly what one would expect given Corollary \ref{cor:deviations-of-prime-divisors-on-a-loglog-scale}.)
The divisors of $p_1 \cdots p_j$ should then lie in a interval of length
\begin{equation*}
    \ll \log p_j \ll 2^{j} e^{-c \sqrt{\log \log y}},
\end{equation*}
and the divisors of $a$ lie in $2^{k - j}$ translates of this interval.
We are thus led to expect
\begin{equation*}
    \prob{\tau(a, y, 2y) \ge 1} \approx \frac{2^k}{\log y} e^{-c \sqrt{\log \log y}},
\end{equation*}
which is much smaller than our original estimate.

This heuristic leads to two key insights in Ford's work: 
\begin{enumerate}
    \item We need a way to capture how the clustering of divisors affects the chance of having a divisor in $(y, 2y]$. 
    Ford achieves this by introducing a quantity $L(a)$ (see Definition \ref{def:L-a} below) that measures the extent of the clustering. 

    \item The main contribution to $H(x, y, 2y)$ should come from integers with $\approx k_0$ prime divisors in $[0, 2y]$ for which the divisors are in fact evenly distributed on a logarithmic scale.
    This is the case when 
    \begin{equation} \label{eq:heuristic-condition-on-prime-factors}
        \log \log p_j \ge j \log 2 - O(1), \qquad 1\le j \le k.
    \end{equation}
    If $\xi_1 \le \dots \le \xi_k$ are uniform order statistics on $[0, 1]$, \eqref{eq:heuristic-condition-on-prime-factors} corresponds to asking that 
    \begin{equation} \label{eq:heuristic-condition-on-order-statistics}
        \xi_j \ge \frac{j - O(1)}{k}, \qquad 1\le j \le k.
    \end{equation}
\end{enumerate}
Ford showed that the probability of \eqref{eq:heuristic-condition-on-order-statistics} occurring is $\asymp 1/k \asymp 1/\log \log y$.
This leads to the new estimate
\begin{equation*}
    H(x, y, 2y) \asymp \frac{x}{(\log y)^\delta (\log \log y)^{3/2}},
\end{equation*}
which finally gives the correct order of magnitude.

\subsection{A global-to-local principle}

To measure the clustering of divisors, we consider the following quantity. 
\begin{definition} \label{def:L-a}
    For any integer $a$, define
    \begin{align*}
        \mcL (a) &= \bigcup \limits_{d\mid a} \left[\log d - \log 2, \log d\right), \\
        L(a) &= \meas(\mcL(a)),
    \end{align*}
    where ``$\meas$'' denotes the standard Lebesgue measure on $\mR$.
\end{definition}

Roughly speaking, $L(a)$ is small when the divisors of $a$ are clustered, and so their corresponding intervals overlap.
Note that an integer $a$ has a divisor in $(y, 2y]$ if and only if 
\begin{equation*}
    \log d \in (\log y, \log y + \log 2]
\end{equation*}
for some $d \mid a$, which is equivalent to $\log y \in \mcL (a)$.

Let us record some easy properties of $L(a)$.
\begin{lemma} \label{lem:properties-of-L}
    We have 
    \begin{enumerate}[label=(\roman*)]
        \item \label{item:properties-1} $L(a) \le \min \Set{\tau (a) \log 2,\, \log a + \log 2}$
        \item \label{item:properties-2} If $(a, b) = 1$, then $L(ab) \le \tau(b) L(a)$.
        \item \label{item:properties-3} If $p_1 < \dots < p_k$ are distinct primes, then 
            \begin{equation*}
                L(p_1 \cdots p_k) \le \min \limits_{0 \le j \le k} 2^{k - j} \log(2p_1 \cdots p_j).
            \end{equation*}
    \end{enumerate}
\end{lemma}

The bound $L(a) \le \tau(a) \log 2$ is useful when $a$ has few prime factors, and so $\tau(a)$ is small.
Item (iii) captures the idea that if one of the prime factors of $a$ is significantly smaller than expected, the divisors of $a$ form clusters, which reduces the size of $L(a)$.

\begin{proof}
    Part (i) follows from the fact that $\mcL(a)$ is the union of $\tau(a)$ (not necessarily disjoint) intervals of length $\log 2$, all contained in $[-\log 2, \log a)$.

    For (ii), note that every divisor $d\mid ab$ can be uniquely written as $d = d_1 d_2$ with $d_1 \mid a$ and $d_2 \mid b$. 
    So
    \begin{align*}
        \mcL (ab) &= \bigcup \limits_{d_2 \mid b} \bigcup \limits_{d_1 \mid a} \left[ \log d_2 + \log d_1 - \log 2, \log d_2 + \log d_1 \right)  \\
            &= \bigcup \limits_{d_2 \mid b} \set{ \log d_2 + t}{t \in \mcL (a)},
    \end{align*}
    which implies $L(ab) \le \tau (b) L(a)$.

    Part (iii) now follows from applying (i) and (ii) with $a = p_1 \cdots p_j$ and $b = p_{j + 1} \cdots p_k$ and noting that $\tau(b) = 2^{k - j}$.
\end{proof}

\begin{remark}
    If $d \mid a$ for squarefree $a$, then
    \begin{equation*}
        \log d = \sum \limits_{p \mid d} \log p.
    \end{equation*}
    Looking at divisors on a logarithmic scale, one can therefore think about the question if an integer has a divisor in $(y, 2y]$ as a subset-sum problem.
    More generally, for $x_1, \dots, x_k \ge 0$, one could define a quantity
    \begin{equation*}
        \mcL(\rho, \mbx) = \bigcup\limits_{\eps_1, \dots, \eps_k \in \Set{0,1}} \left[ -\rho + \sum \limits_{i = 1}^k \eps_i x_i, \sum \limits_{i = 1}^k \eps_i x_i \right)
    \end{equation*}
    to measure the distribution of the subset sums of $x_1, \dots, x_k$.
\end{remark}

Suppose $n \le x$ is chosen uniformly at random and write $n = ab$ with $P^+(a) \le 2y < P^-(b)$ as before.
Then $H(x, y, 2y)$ is equal to the number of $n \le x$ for which $\log y \in \mcL (a)$.
Instead of trying to compute the probability that $a$ has a divisor in $(y, 2y]$ directly, the problem can be reduced to estimating the size of a certain weighted average \mbox{of $L(a)$}.
Ford calls this a ``global-to-local'' principle, as $L(a)$ captures the global distribution of the divisors \mbox{of $a$.}

Assuming again that $\log a \asymp \log y$, we expect a positive proportion of $\mcL(a)$ to be of order $\asymp \log y$, and so the probability that $\tau(n, y, 2y) \ge 1$ should be roughly $\asymp \expn{L(a)} / \log y$.
(Here, $a$ should be thought of as a function of $n$.)
Recall from Section \ref{subsec:kubilius-model} that when $y$ is small compared to $x$, the exponents in the prime factorisation of $a$ are approximately geometrically distributed. 
Assuming this is still reasonable for the range we are interested in, we obtain that the probability of choosing one particular $a$ is roughly
\begin{equation*}
    \approx \frac{1}{a} \cdot \prod\limits_{p \le y} \left(1 - \frac{1}{p}\right) \asymp \frac{1}{\log y} \cdot \frac{1}{a}.
\end{equation*}
This leads us to expect that
\begin{equation*}
    \prob{\tau(n, y, 2y) \ge 1} \asymp \frac{1}{\log y} \expn{L(a)} \asymp \frac{1}{(\log y)^2} \sum \limits_{a} \frac{L(a)}{a}.
\end{equation*}
Such a statement holds indeed; this is made precise in Lemma \ref{lem:global-to-local-principle}.

Define
\begin{equation*}
    \mcS(y) = \set{ n \in \mN}{n \text{ squarefree},  \; P^+(n) \le y},
\end{equation*}
and note that this is a finite set.
Furthermore, let $y_0$ be a large constant.

\begin{lemma}[Global-to-local principle]
    \label{lem:global-to-local-principle}
    For $y \ge y_0$, we have
    \begin{equation*}
        H(x, y, 2y) \asymp \frac{x}{(\log y)^2} \sum \limits_{a \in \mcS(y)} \frac{L(a)}{a}.
    \end{equation*}    
\end{lemma}

The lemma is one of the key steps in the proof of Theorem \ref{thm:integers-with-a-divisor-in-y-2y}, serving as a starting point for both the lower and the upper bounds.
The remainder of this section is devoted to the proof of the lemma.
We will provide the details for the lower bound and outline the proof of the upper bound, which is more technically involved. 

Let us start by collecting two estimates showing that for sums over $L(a)/a$, we have some flexibility in choosing the set of integers $a$ summed over.
In particular, when summing over $a \in \mcS(y_1)$, taking $y_2 = y_1^{O(1)}$ changes the value of the sum by at most a constant factor.

\begin{lemma} \label{lem:replace-sum-over-a}
    \begin{enumerate}
        \item Suppose $2 \le y_1 < y_2$. Then 
            \begin{equation*}
                \sum \limits_{a \in \mcS(y_2)} \frac{L(a)}{a} \ll \left( \frac{\log y_2}{\log y_1} \right)^2 \sum \limits_{a \in \mcS(y_1)}\frac{L(a)}{a}.
            \end{equation*}
        \item For any $y \ge 2$,
            \begin{equation*}
                \sum \limits_{a \le y} \frac{L(a)}{a} \asymp \sum \limits_{a \in \mcS(y)} \frac{L(a)}{a}.
            \end{equation*}
    \end{enumerate}
\end{lemma}

\begin{proof}
    (i) Write every $a \in \mcS(y_2)$ uniquely as $a = a_1 a_2$ with $P^+(a_1) \le y_1 < P^- (a_2)$.
    Then $(a_1, a_2 ) = 1$, so Lemma \ref{lem:properties-of-L} (ii) implies that $L(a) \le \tau(a_2) L(a_1)$.
    Hence,
    \begin{equation*}
        \sum \limits_{a \in \mcS(y_2)} \frac{L(a)}{a} \le \sum \limits_{a_1 \in \mcS(y_1)} \frac{L(a_1)}{a_1} \sum \limits_{\substack{a_2 \in \mcS(y_2) \\ P^-(a) > y_1}} \frac{\tau(a_2)}{a_2}.
    \end{equation*}
    The inner sum on the right-hand side is
    \begin{align*}
        \prod \limits_{y_1 < p \le y_2}  \left(1 + \frac{2}{p} \right) \le \exp \left( 2 \sum \limits_{y_1 < p \le y_2} \frac{1}{p} \right) \ll \left(\frac{\log y_2}{\log y_1}\right)^2
    \end{align*}
    by Mertens' second estimate. 

    (ii) The ``$\ll$'' part is clear.
    For the other direction, note the inequality
    \begin{equation*}
        \ind{a \le y} \ge 1 - \frac{\log a}{\log y}.
    \end{equation*}
    We use the above to write
    \begin{align} \label{eq:lower-bound-ineq}
        \sum \limits_{a \le y} \frac{L(a)}{a} &\ge \sum \limits_{a \in \mcS(y^{1/4})} \ind{a \le y} \cdot \frac{L(a)}{a}\nonumber \\&\ge \sum \limits_{a \in \mcS(y^{1/4})} \frac{L(a)}{a}  - \frac{1}{\log y} \sum \limits_{a \in \mcS(y^{1/4})} \frac{L(a) \log a}{a}.
    \end{align}
    Now
    \begin{align*}
        \sum \limits_{a \in \mcS(y^{1/4})} \frac{L(a) \log a}{a} 
            &= \sum \limits_{a \in \mcS(y^{1/4})} \frac{L(a)}{a} \sum\limits_{p \mid a} \log p \\
            &= \sum \limits_{p \le y^{1/4}} \frac{\log p}{p} \sum\limits_{\substack{b \in \mcS(y^{1/4})\\p \nmid b}} \frac{L(pb)}{b} \\
            &\le 2\sum \limits_{p \le y^{1/4}} \frac{\log p}{p} \sum\limits_{b \in \mcS(y^{1/4})} \frac{L(b)}{b} \\
            &= \left({\tfrac{1}{2}} \log y + O(1) \right) \sum\limits_{b \in \mcS(y^{1/4})} \frac{L(b)}{b},
    \end{align*}
    where we have used Lemma \ref{lem:properties-of-L} (ii) for the third inequality and Mertens' first estimate to pass to the last line.
    Combining this with \eqref{eq:lower-bound-ineq}, we get
    \begin{equation*}
        \sum \limits_{a \le y} \frac{L(a)}{a} \ge \left(1 - \frac{1}{2} - O\left( \frac{1}{\log y} \right) \right)\sum \limits_{a \in \mcS(y^{1/4})} \frac{L(a)}{a} \gg \sum \limits_{a \in \mcS(y^{1/4})} \frac{L(a)}{a}.
    \end{equation*}
    Finally, part (i) tells us that we can replace $\mcS(y^{1/4})$ by $\mcS(y)$.
\end{proof}

We shall also use the following simple sieve bound.
\begin{lemma} \label{lem:sieve-bound}
    Let $z \ge 2$ and $x \ge z$. Then
    \begin{equation*}
        \#\set{n \le x}{P^-(n) > z} \ll \frac{x}{\log z}.
    \end{equation*}
    If $x \ge 2z$, we also have
    \begin{equation*}
        \#\set{n \le x}{P^-(n) > z} \gg \frac{x}{\log z}.
    \end{equation*}
\end{lemma}
Heuristically, the lemma seems plausible if one assumes that divisibility by distinct primes is approximately independent; the number of integers $\le x$ not divisible by any prime $p \le z$ should then be
\begin{equation*}
    \approx x \prod_{p \le z} \left( 1 - \frac{1}{p} \right),
\end{equation*}
which is $\asymp x/ \log z$ by Mertens' third estimate.
A proof of Lemma \ref{lem:sieve-bound} is sketched in \cite[Exercise 02]{divisors}.

We now have all the tools to prove the lower bound in Lemma \ref{lem:global-to-local-principle}.
The idea is to count integers $n \le x$ of the form $n = apb$, where $p$ is a prime in $(y^{1 - \delta}, 2y]$, for some $\delta > 0$, and $ap$ has a divisor in $(y, 2y]$. 
Any such integer $n$ is counted by $H(x, y, 2y)$, so a lower bound on the number of choices for $n$ will give a lower bound on $H(x, y, 2y)$.
\begin{proof}[Proof of Lemma \ref{lem:global-to-local-principle}, lower bound]
    Consider integers $n \le x$ of the form $n = apb$, where $a \le y^{1/8}$, $p$ is prime, and all prime factors of $b$ are $> 2y$ or in $(y^{1/4}, y^{3/4}]$.
    Moreover, assume that $a$ has a divisor $d$ for which $dp \in (y, 2y]$.
    The last condition is equivalent to $\log(y/p) \in \mcL(a)$, and it implies that $n$ is counted by $H(x, y, 2y)$. 
    Furthermore, it implies that $p \in (y^{7/8}, 2y]$.
    This means that the representation of $n$ as $abp$ is unique (if it exists), since $p$ and the prime factors of $a$ and $b$ lie in disjoint intervals.

    For a fixed pair $a, p$, we now estimate the number of choices for $b$.
    Note that we must have $b \le x/ap$.
    Also, $x \ge y^2$ and $ap \le 2y^{9/8}$, so $x/ap \ge \frac{1}{2}y^{7/8}$.
    When $x / ap \ge 4y$, the sieve bound from Lemma \ref{lem:sieve-bound} implies that 
    \begin{equation*}
        \#\set{b \le \frac{x}{ap}}{P^-(b) > 2y} \gg \frac{x}{ap \log y}.
    \end{equation*}
    When $x/ap < 4y$, we need to use that $b$ can have prime factors in $(y^{1/4}, y^{3/4}]$.
    We claim that there are $\gg x/(ap \log y)$ possibilities to choose $b$ as a product of two primes $q_1, q_2$ from that interval.
    Indeed, suppose that we choose the first prime $q_1 \in (4y^{1/4}, y^{1/2}]$.    
    Then 
    \begin{equation*}
         \frac{x}{apq_1} \le \frac{4y}{q_1} \le y^{3/4},
    \end{equation*}
    so we can choose $q_2$ from the interval $(y^{1/4}, x/(apq_1)]$.
    Since $x/(apq_1) \gg y^{7/8} / q_1 \gg y^{3/8}$,
    the PNT tells us that this interval contains 
    \begin{equation*}
        \pi\left( \frac{x}{apq_1} \right) - \pi\left(y^{1/4}\right)\gg \frac{x}{apq_1 \log y} - y^{1/4} \gg \frac{x}{apq_1 \log y}
    \end{equation*}
    primes. 
    Summing over all choices for $q_1$ and applying Mertens' second estimate \eqref{eq:mertens-second}, we conclude that there are
    \begin{equation*}
        \gg \frac{x}{a p \log y} \underbrace{\sum \limits_{4y^{1/4} < q_1 \le y^{1/2}} \frac{1}{q_1}}_{\asymp 1} \gg \frac{x}{a p \log y}
    \end{equation*}
    choices for $b$.

    Together, the number of admissible $n$ is
    \begin{equation} \label{eq:lower-bound-choices-for-n}
        \gg \sum \limits_{a \le y^{1/8}} \frac{x}{a \log y} \sum \limits_{p \colon \log (y/p) \in \mcL(a)} \frac{1}{p}. 
    \end{equation}
    To estimate the sum over $p$, we apply Mertens' estimate to every interval in $\mcL(a)$. 
    This is encapsulated in the following claim.
    \begin{claim*}
        \begin{equation*}
            \sum \limits_{p \colon \log (y/p) \in \mcL(a)} \frac{1}{p} \asymp \frac{L(a)}{\log y}.
        \end{equation*}
    \end{claim*}
    \begin{proof}
        \renewcommand{\qedsymbol}{$\Diamond$}
        Recall that $\mcL(a)$ is the union of intervals of length at least $\log 2$, all of which are contained in $[0, \frac{1}{8}\log y)$ since $a \le y^{1/8}$.
        For every such interval $[u, v)$, we have
        \begin{equation*}
            \log\left(\frac{y}{p} \right) \in [u, v) \quad \Longleftrightarrow \quad p \in (ye^{-v}, ye^{-u}].
        \end{equation*}
        Mertens' second estimate \eqref{eq:mertens-second} (with the strong error term) gives
        \begin{align*}
            \sum \limits_{u \le \log (y/p) < v} \frac{1}{p} 
            % &=\sum \limits_{ye^{-v} < p \le ye^{-u}} \frac{1}{p} \\ 
            &= \log \log (ye^{-u}) - \log \log (ye^{-v}) + O\left(e^{-c\sqrt{-v+\log y}}\right) \\
            % &= \int_{ye^{-v}}^{ye^{-u}} \frac{1}{t \log t}\,\mathrm d t + O\left(e^{-c\sqrt{-v+\log y}}\right) \\ 
            &= \int_{\log y -v}^{\log y -u} \frac{1}{t}\,\mathrm d t + O\left(e^{-c\sqrt{\log y}}\right) \\ 
            &\asymp \frac{v - u}{\log y} + O\left(e^{-c\sqrt{\log y}}\right).
        \end{align*}
        For the last step, we used that the integrand in the second line is of order $1/\log y$ throughout, since $u, v\le \frac{1}{8} \log y$.
        Since $v - u \ge \log 2$, we can omit the error term. 
        The claim now follows by summing over all intervals $[u, v)$.
    \end{proof}

    Combining the claim and \eqref{eq:lower-bound-choices-for-n}, we arrive at
    \begin{equation*}
        H(x, y, 2y) \gg \frac{x}{(\log y)^2} \sum \limits_{a \le y^{1/8}} \frac{L(a)}{a}.
    \end{equation*}
    Finally, parts (i) and (ii) of Lemma \ref{lem:replace-sum-over-a} tell us that we can replace the sum over $a \le y^{1/8}$ by a sum over $\mcS(y)$, which concludes the proof.
\end{proof}

\begin{remark}
    For the purpose of proving Theorem \ref{thm:integers-with-a-divisor-in-y-2y}, replacing the sum over $a \le y^{1/8}$ by $a \in \mcS(y)$ is not necessary, and this is not done in \cite{integerswithadivisor}. 
    However, it will slightly simplify later calculations; and it gives a more symmetric lemma, with a lower bound of the same shape as the upper bound.
\end{remark}

In the remainder of this section, we briefly sketch the derivation of the upper bound in Lemma \ref{lem:global-to-local-principle}.
The first step is to relate $H(x, y, 2y)$ to the number of \textit{squarefree} integers with a divisor in a given interval.
Let $H^*(x, y, 2y)$ denote the number of squarefree integers $\le x$ with a divisor in $(y, 2y]$.
By decomposing every integer $\le x$ into coprime squarefree and squarefull parts, one quickly sees that
\begin{equation} \label{eq:relate-H-x-y-2y-to-squarefree-integers}
    H(x, y, 2y) \le \sum\limits_{\substack{n \le y^\delta\\ n \text{ squarefull}}} \sum \limits_{d \mid n} H^*\left( \frac{x}{n}, \frac{y}{d}, \frac{2y}{d} \right) + O\left( \frac{x}{y^{\delta/2}} \right),
\end{equation}
where the error term comes from omitting integers whose squarefull part is $> y^\delta$.

Next, one relates the quantities $H^*(x / n, y/d, 2y/d)$ to a sum involving $L(a)$ by proving that for large enough $y \le x^{0.51}$,
\begin{equation} \label{eq:sum-over-L-a-with-weird-denominator}
    H^*(x, y, 2y) - H^{*}({\textstyle \frac{1}{2}}x, y, 2y) \ll x \sum \limits_{a \in \mcS(2y)} \frac{L(a)}{a \log^2 \max\left(y^{3/4}/a,\; P^+(a)\right)}.
\end{equation}
The strategy here is similar to the proof of the lower bound of Lemma \ref{lem:global-to-local-principle}. One writes each squarefree $n \in (x/2, x]$ with a divisor in $(y, 2y]$ as $n = apb$ with $P^+(a) < p < P^-(b)$ in a strategic way, and then uses the sieve bound from Lemma \ref{lem:sieve-bound} and Mertens' estimate to count the number of admissible $n$.

One can then apply the above to each summand $H^*(x / n, y/d, 2y/d)$ in \eqref{eq:relate-H-x-y-2y-to-squarefree-integers} using dyadic summation. 
Finally, to pass from \eqref{eq:sum-over-L-a-with-weird-denominator} to a sum over $L(a)/a$ of the desired form, one can apply a technical lemma due to Koukoulopoulos \cite[Lemma 2.2 (b)]{localizedfactorizationsofintegers} (see also \cite{integerswithadivisor} for a simple proof).
The proof of the upper bound can be concluded by summing over pairs $d \mid n$ and noting that 
\begin{equation*}
    \sum \limits_{n \text{ squarefull}} \frac{\tau(n)}{n} = \prod\limits_{p} \left(1 + \frac{3}{p^2} + \frac{4}{p^3} + \dots  \right) \ll 1.
\end{equation*}

For details, we refer the reader to \cite[Section 7.6]{fordlecturenotes}.

\subsection{Proof of the lower bound}
    \label{subsec:proof-of-lower-bound}

In the preceding section (Lemma \ref{lem:global-to-local-principle}), we showed that 
\begin{equation*}
    H(x, y, 2y) \gg \frac{x}{(\log y)^2} \sum \limits_{a \in \mcS(y)} \frac{L(a)}{a}
\end{equation*}
for $y \ge y_0$.
To conclude the proof of the lower bound in Theorem \ref{thm:integers-with-a-divisor-in-y-2y}, we now want to further bound the right-hand side.
We reduce the problem of estimating the weighted average over $L(a)$ to a more combinatorial problem, namely, counting pairs of divisors that are close to each other. 
Define
\begin{equation*}
    W(a) = \# \set{(d, d')}{d,d' \mid a,\, \left| \log \frac{d}{d'}\right| \le \log 2}.
\end{equation*}
Note that a pair of divisors $(d, d')$ is counted by $W(a)$ whenever their respective intervals in $\mcL(a)$ overlap, so $W(a)$ is large whenever the divisors of $a$ are clustered on a logarithmic scale.

The next lemma essentially plays the role of a second-moment bound.

\begin{lemma} \label{lem:snd-moment-bound}
    For any finite set $\mcA \In \mN$, we have
    \begin{equation*}
        \sum \limits_{a \in \mcA} \frac{L(a)}{a} \ge (\log 2) \frac{\left(\sum_{a \in \mcA} \frac{\tau (a)}{a} \right)^2}{\sum_{a \in \mcA} \frac{W(a)}{a}}.
    \end{equation*}
\end{lemma}

\begin{proof}
    First, observe that 
    \begin{equation*}
        \int \tau(a, e^t, 2 e^t)\, \mathrm{d} t = (\log 2) \tau(a),
    \end{equation*}
    since every divisor is counted in an interval of length $\log 2$ by the integral. 
    By the Cauchy--Schwarz inequality, 
    \begin{align*}
        \left(\int \tau(a, e^t, 2 e^t)\, \mathrm{d} t\right)^2 &= \left(\int \ind{t \in \mcL(a)} \tau(a, e^t, 2 e^t)\, \mathrm{d} t \right)^2 \\
            &\le \int \ind{t \in \mcL(a)}\, \mathrm{d}t \cdot  \int \tau(a, e^t, 2 e^t)^2\, \mathrm{d} t
            = L(a) \int \tau(a, e^t, 2 e^t)^2\, \mathrm{d} t.
    \end{align*}
    A second application of Cauchy--Schwarz (to the sum over $a$) gives
    \begin{align*}
        (\log 2)^2\left( \sum \limits_{a \in \mcA} \frac{\tau(a)}{a} \right)^2 
            &= \left( \sum \limits_{a \in \mcA} \frac{1}{a} \int \tau(a, e^t, 2 e^t)\, \mathrm{d} t \right)^2 \\
            &\le \left( \sum \limits_{a \in \mcA} \frac{L(a)}{a} \right) \left( \sum \limits_{a \in \mcA}\frac{1}{a} \int \tau(a, e^t, 2 e^t)^2\, \mathrm{d} t \right).
    \end{align*}
    Now
    \begin{align*}
        \int \tau(a, e^t, 2 e^t)^2\, \mathrm{d} t 
            &= \sum \limits_{d, d' \mid a} \int \ind{d, d' \in (e^t, 2e^t]}\, \mathrm{d}t \\
            &= \sum \limits_{d, d' \mid a} \max \Set{0, \log 2 - \left| \frac{\log d}{\log d'} \right|} 
            \le (\log 2) W(a).
    \end{align*}
    The result follows after rearranging.
\end{proof}

We want to apply the above lemma to sets of integers whose prime factors are localised, in the sense that the number of prime factors that lie in a certain interval is prescribed. 
To this end, we partition the primes into sets $(D_j)_{j \ge 1}$, defined as follows. 
Set $\lambda_0 = 3/2$, and for $j \ge 1$, define $\lambda_j$ inductively as the largest prime such that
\begin{equation} \label{eq:definition-of-lambda-j}
    \sum \limits_{\lambda_{j - 1} < p \le \lambda_j} \frac{1}{p} \le \log 2.
\end{equation}
Now let $D_j = (\lambda_{j - 1}, \lambda_j]$.
Furthermore, for any vector $\mbb = (b_1, \dots, b_J)$ whose entries are nonnegative integers, let $\mcA (\mbb)$ denote the set of squarefree integers with precisely $b_j$ prime factors in $D_j$ for $1 \le j \le J$, and no other prime factors.

The sequence $(\lambda_j)_{j \ge 0}$ grows extremely rapidly. Roughly speaking, the following observation shows that $\log \lambda_j \approx 2 \log \lambda_{j - 1}$, and so $\lambda_j \approx \lambda_{j - 1}^2$.
\begin{observation}
    We have $\log \lambda_j = 2^{j + O(1)}$.
\end{observation}
\begin{proof}
    Observe that since $\lambda_j$ is chosen to be the largest prime such that \eqref{eq:definition-of-lambda-j} holds, the left-hand side in \eqref{eq:definition-of-lambda-j} is $\log 2 + O(1/\lambda_j)$.
    Now Mertens' estimate \eqref{eq:mertens-second} implies that
    \begin{equation} \label{eq:mertens-applied-to-D-j}
        \log \log \lambda_j - \log \log \lambda_{j - 1} = \log 2+ O\left( e^{-c \sqrt{\log \lambda_{j - 1}}} \right).
    \end{equation}
    Applying this iteratively, we obtain
    \begin{equation} \label{eq:mertens-applied-iteratively-to-D-j}
        \log \log \lambda_j = j \log 2 + O\left(\sum \limits_{i = 0}^{j - 1} e^{-c \sqrt{\log \lambda_{i}}} \right),
    \end{equation}
    and the error term is $O(1)$. 
    To see this, observe that for large enough $j$, the error term in \eqref{eq:mertens-applied-to-D-j} is $< 0.1$ (say). So $\log \log \lambda_{j} - \log \log \lambda_{j - 1} \ge 1/2$, which implies $\log \lambda_j \ge \sqrt{e} \log \lambda_{j - 1}$.
    Hence, the sequence $\log \lambda_{j}$ grows exponentially in $j$, and so the summands in the big-O term decay very fast.
    Finally, exponentiating \eqref{eq:mertens-applied-iteratively-to-D-j}, we conclude that $\log \lambda_j  = 2^{j + O(1)}$.
\end{proof}

The intuition behind the above construction is that, when summing over reciprocals of primes, every interval $D_j$ carries roughly the same weight (namely $\approx \log 2$). 
As a general principle, 
\begin{equation*}
    \sum \limits_{a \in \mcA (\mbb)} \frac{1}{a} \le \frac{1}{b_1! \cdots b_J!} \prod \limits_{j = 1}^J \left( \sum \limits_{p \in D_j} \frac{1}{p} \right)^{b_j} \le \frac{(\log 2)^{b_1 + \dots + b_J}}{b_1 ! \cdots b_J!},
\end{equation*}
since each of the sums over $p \in D_j$ is $\le \log 2$.
Under some mild condition on the numbers $b_j$, a similar inequality holds in the other direction; see Lemma \ref{lem:lower-bound-for-sum-over-tau-a}.

Also recall from Section \ref{subsec:kubilius-model} that heuristically, the prime divisors of a random integer behave roughly like a Poisson point process with intensity 1 on a log-log-scale.
In light of this, we expect a ``typical'' integer $a$ to have $\Theta(1)$ prime factors in each $D_j$, at least when $j$ is neither too small nor too large.

In what follows, let $y_0$ and $M$ be some large (but fixed) constants and assume that $y \ge y_0$.
Put
\begin{equation*}
    k = \lflr \frac{\log \log y}{\log 2} - 2M \rflr,
\end{equation*}
and 
\begin{equation*}
    J = M + k - 1.
\end{equation*}
We now restrict our attention to integers with a precise count of prime factors in $D_M, \dots, D_J$ and $k$ prime factors in total.
Set
\begin{equation*}
    \mcB = \set{(b_1, \dots, b_J)}{b_j = 0 \text{ for } j < M,\, b_1 + \dots + b_J = k}.
\end{equation*}
One should think about the entries $b_j$ of a vector in $\mcB$ as the number of prime divisors a given integer has in $D_j$.
The choice of $J$ ensures that on average, an integer corresponding to such a vector has one prime divisor in each of $D_M, \dots, D_J$. 
Furthermore, for $\mbb \in \mcB$ and $a \in \mcA(\mbb)$, we have $a \in \mcS(y)$.
Indeed, in that case,
\begin{equation*}
    \log P^+(a) \le \log \lambda_J \le 2^{J + O(1)} \le 2^{\log \log y / \log 2} = \log y,
\end{equation*}
provided that $M$ is large enough to absorb the $O(1)$ term.
Hence, we can write
\begin{equation} \label{eq:reduce-to-b-in-B}
    \sum \limits_{a \in \mcS(y)} \frac{L(a)}{a} \ge \sum \limits_{\mbb \in \mcB} \sum \limits_{a \in \mcA(\mbb)} \frac{L(a)}{a}.
\end{equation}
At first sight, this might seem somewhat wasteful. 
However, recall that, heuristically, the main contribution to $H(x, y, 2y)$ should come from integers with $\log \log y / \log 2 + O(1)$ prime factors $\le 2y$.
Furthermore, using Lemma \ref{lem:properties-of-L} (ii), one can easily show that by restricting the sum to numbers $a$ without small prime factors, one only loses a constant factor. 
Therefore, we may still hope to recover the correct order of magnitude from the right-hand side of \eqref{eq:reduce-to-b-in-B}.

Motivated by Lemma \ref{lem:snd-moment-bound} and the above discussion, we prove the next two lemmas.
We further restrict the range of $\mbb$ to avoid numbers whose prime divisors are distributed too irregularly. 
To this end, define
\begin{equation*}
    \mcB^* = \set{\mbb \in \mcB}{b_j \le Mj \text{ for } j \ge M}.
\end{equation*}
Note that $b_j \le Mj$ is a rather mild constraint, as typically we expect only $\sim 1$ prime in every interval $D_j$.

\begin{lemma} \label{lem:lower-bound-for-sum-over-tau-a}
    For $\mbb \in \mcB^*$, we have
    \begin{equation*}
        \sum \limits_{a \in \mcA(\mbb)} \frac{\tau (a)}{a} \gg \frac{(2 \log2 )^k}{b_M! \cdots b_J!}.
    \end{equation*}
\end{lemma}

\begin{remark}
    Without any growth assumption on the entries of $\mbb$, the above is not true. 
    For example, the vector $\mbb$ with $b_M = k$ and $b_j = 0$ for $j \ne M$ lies in $\mcB$, but $\mcA(\mbb) = \emptyset$ if $y$ is too big, since there is only a bounded number of distinct primes in $D_M$.
\end{remark}

\begin{proof}[Proof of Lemma \ref{lem:lower-bound-for-sum-over-tau-a}]
    Note that $\tau(a) = 2^k$ for all $a \in \mcA(\mbb)$.
    Hence, 
    \begin{align*}
        \sum \limits_{a \in \mcA(\mbb)} \frac{\tau (a)}{a} &= 2^k \prod \limits_{j = M}^J \frac{1}{b_j!} \left( \sum \limits_{p_1 \in D_j} \frac{1}{p_1} \sum \limits_{p_2 \in D_j \setminus \Set{p_1}} \frac{1}{p_2} \cdots \sum \limits_{p_{b_j} \in D_j \setminus \{p_1, \dots, p_{b_j -1}\}} \frac{1}{p_{b_j}}\right) \\
        &\ge 2^k \prod \limits_{j = M}^J \frac{1}{b_j!} \left(\log 2 - \frac{b_j}{\lambda_{j - 1}} \right)^{b_j} \\ 
        &\ge \frac{(2 \log 2)^k}{2 b_M! \cdots b_J!}.
    \end{align*}
    For the last line, we used that
    \begin{equation} \label{eq:bound-on-product}
        \prod\limits_{j = M}^J\left(1 - \frac{b_j}{\lambda_{j - 1} \log 2}\right)^{b_j} \ge \frac{1}{2},
    \end{equation}
    provided that $M$ is large enough.
    To see this, observe that for $0 \le x \le 1/2$, we have the inequality $1 - x \ge e^{-2x}$. 
    Hence, using that $b_j \le Mj$ for all $j$,
    \begin{equation*}
        \prod\limits_{j = M}^J\left(1 - \frac{b_j}{\lambda_{j - 1} \log 2}\right)^{b_j} \ge \exp\left(-\sum_{j = M}^J \frac{2b_j^2}{\lambda_{j - 1} \log 2} \right) \ge \exp\left(-\sum_{j = M}^J \frac{2(Mj)^2}{\lambda_{j - 1} \log 2} \right).
    \end{equation*}
    Since $\lambda_{j - 1} = \exp(2^{j + O(1)})$, the sum in the exponent is rapidly convergent and in particular, it is $< \log 2$ if $M$ is large enough.
    This shows \eqref{eq:bound-on-product} and thus completes the proof.
\end{proof}

\begin{lemma} \label{lem:upper-bound-for-sum-over-W-a}
    For any $\mbb = (b_1, \dots, b_J) \in \mcB$, we have
    \begin{equation*}
        \sum \limits_{a \in \mcA (\mbb)} \frac{W(a)}{a} \ll \frac{(2 \log 2)^{k}}{b_M! \cdots b_J!} f(\mbb),
    \end{equation*}
    where 
    \begin{equation*}
        f(\mbb) = \sum \limits_{j = M}^J 2^{(b_M - 1) + \dots + (b_j - 1)}.
    \end{equation*}
\end{lemma}

\begin{proof}
    For $M\le j \le J$, put $B_j = b_M + \dots + b_j$ (note that $b_j = 0$ for $j < M$).
    Consider $a \in \mcA (\mbb)$, and write $a = p_1 \cdots p_k$, with
    \begin{equation} \label{eq:localised-prime-factors}
        p_{B_{j - 1} + 1}, \dots, p_{B_j} \in D_j, \quad M \le j \le J.
    \end{equation}
    There is a 1-to-1 correspondence between the divisors of $a$ and subsets $X \In \Set{1, \dots, k}$. 
    So recalling the definition of $W(a)$, we see that $W(a)$ is precisely equal to the number of pairs $X, Y \In \Set{1, \dots, k}$ for which 
    \begin{equation} \label{eq:equation-for-W(a)}
        \left| \sum \limits_{i \in X} \log p_i - \sum \limits_{i \in Y} \log p_i \right| \le \log 2.
    \end{equation}
    This implies that 
    \begin{equation} \label{eq:rewrite-in-terms-of-X-and-Y}
        \sum \limits_{a \in \mcA(\mbb)} \frac{W(a)}{a} \le \frac{1}{b_M ! \cdots b_J!} \sum \limits_{X, Y \In \Set{1, \dots, k}} \sum \limits_{\substack{p_1, \dots, p_k\\ \text{satisfying \eqref{eq:localised-prime-factors} and \eqref{eq:equation-for-W(a)}}  }} \frac{1}{p_1 \cdots p_k},
    \end{equation}
    where the factor $1/(b_M ! \cdots b_J!)$ comes from the fact that for each $a \in \mcA(\mbb)$ we can choose the primes $p_{B_{j - 1} + 1}, \dots, p_{B_j} \in D_j$ in $b_j!$ different orders. Also note that this really is an inequality since the primes in the inner sum need no longer be distinct.

    First, consider the contribution of pairs $X = Y$. There are precisely $2^k$ such pairs, and for each of them, the innermost sum in \eqref{eq:rewrite-in-terms-of-X-and-Y} is
    \begin{equation*}
        \left( \sum \limits_{p \in D_M} \frac{1}{p} \right)^{b_M} \cdots \left( \sum \limits_{p \in D_J} \frac{1}{p} \right)^{b_J} \le (\log 2) ^k.
    \end{equation*}

    Now consider $X \ne Y$. 
    For fixed $I \in \Set{1, \dots, k}$, we want to estimate the contribution of pairs $X, Y$ with $I = \max (X \triangle Y)$, i.e., for which $I$ is the largest element in which $X$ and $Y$ differ.
    Let $e(I) \in \Set{M, \dots, J}$ be such that $B_{e(I) - 1} < I \le B_{e(I)}$.
    Suppose that all $p_i$ are fixed except for $p_I$ (which varies over $D_{e(I)}$).
    Then \eqref{eq:equation-for-W(a)} implies that $p_I$ is restricted to an interval $[U, 4U]$ for some number $U$ (depending on the other $p_i$), since $\log p_I$ is restricted to an interval of length $2\log 2$.
    We then obtain from Mertens' estimate \eqref{eq:mertens-second} that 
    \begin{align} \label{eq:bound-using-U}
        \sum \limits_{\substack{p_I \in D_{e(I)}\\U \le p_I \le 4U}} \frac{1}{p_I} 
            &\le \sup \limits_{\lambda_{e(I) - 1} \le u \le \lambda_{e(I)}} \sum \limits_{u < p_I \le 4u} \frac{1}{p_I}\nonumber \\
            &\le \sup \limits_{\lambda_{e(I) - 1} \le u \le \lambda_{e(I)}} (\log \log 4u) - \log \log u + O\left(e^{-c\sqrt{\log u}}\right)\nonumber \\
            &\ll \sup \limits_{\lambda_{e(I) - 1} \le u \le \lambda_{e(I)}} \frac{1}{\log u}\nonumber \\
            &\ll \frac{1}{\log \lambda_{e(I) - 1}}\nonumber \\
            &\ll 2^{-e(I)}.
    \end{align}
    For the last inequality, we used $\log \lambda_{e(I) - 1} \ge 2^{e(I) + O(1)}$.
    Note the role of $U$ in restricting the range of the sum: If we took the sum over all of $D_{e(I)}$, we would only get the trivial upper bound $\log 2$.

    Summing over the $p_i$ with $i \ne I$, it follows from \eqref{eq:bound-using-U} that for a fixed pair $X \ne Y$ with $I = \max (X \triangle Y)$, the inner sum in \eqref{eq:rewrite-in-terms-of-X-and-Y} is at most $(\log 2)^k \cdot 2^{-e(I)}$. 
    There are precisely\footnotemark
    \footnotetext{We have two choices to put $I$ either in $X$ or in $Y$. Next, we have $2^{k-I}$ choices to include each number $> I$ in both $X$ and $Y$ or in neither set. Finally, for the numbers $< I$ we have $4^{I - 1}$ choices, since we can include each number in either of $X$ and $Y$, in both, or in neither of the two.}
    \begin{equation*}
        2 \cdot 2^{k- I} \cdot 4^{I - 1}  = 2^{k + I - 1}
    \end{equation*}
    choices for $X, Y$ with $I = \max (X \triangle Y)$,
    so the total contribution of this particular $I$ is 
    \begin{equation*}
        \ll 2^{k + I - 1 - e(I)} (\log 2) ^k.
    \end{equation*}
    Putting everything together, we obtain
    \begin{align*}
        \sum \limits_{a \in \mcA(\mbb)} \frac{W(a)}{a} 
            &\ll \frac{(2 \log 2)^k}{b_M ! \cdots b_J!} \left( 1 + \sum \limits_{I = 1}^k 2^{I - 1 -e(I)} \right) \\
            &\le \frac{(2 \log 2)^k}{b_M ! \cdots b_J!} \left( 1 + \sum \limits_{j = M}^J 2^{-j} \sum \limits_{B_{j - 1} < I \le B_j} 2^{I - 1} \right) \\ 
            &\le  \frac{(2 \log 2)^k}{b_M ! \cdots b_J!} \left( 1 + \sum \limits_{j = M}^J 2^{B_j -j} \right).
    \end{align*}
    Recalling that $B_j = b_M + \dots + b_j$, the right-hand side is equal to
    \begin{equation*}
            \frac{(2 \log 2)^k}{b_M ! \cdots b_J!}  \left( 1 + 2^{1 - M} \sum \limits_{j = M}^J 2^{(b_M - 1) + \dots + (b_j - 1)} \right)
            \ll \frac{(2 \log 2)^k}{b_M ! \cdots b_J!}  \left( 1 + f(\mbb) \right).
    \end{equation*}
    Finally, since the summand for $j = M$ is at least $2^{-1}$, we have $f(\mbb) \ge 1/2$, and so we may drop the $+1$ term. 
    This concludes the proof.
\end{proof}

Combining Lemmas \ref{lem:snd-moment-bound}, \ref{lem:lower-bound-for-sum-over-tau-a}, and \ref{lem:upper-bound-for-sum-over-W-a}, we have
\begin{align} \label{eq:estimate-h-x-y-2y-by-sum-over-f-b}
    H(x, y, 2y) &\gg \frac{x}{(\log y)^2} \sum \limits_{\mbb \in \mcB^*} \sum \limits_{a \in \mcA(\mbb)} \frac{L(a)}{a} \nonumber \\ 
    &\gg \frac{x}{(\log y)^2} \sum \limits_{\mbb \in \mcB^*} \frac{\left(\sum_{a \in \mcA(\mbb)} \frac{\tau (a)}{a} \right)^2}{\sum_{a \in \mcA(\mbb)} \frac{W(a)}{a}} \nonumber \\
    &\gg \frac{x(2 \log 2)^k}{(\log y)^2}  \sum \limits_{\mbb \in \mcB^*} \frac{1}{b_M! \cdots b_J! f(\mbb)}.
\end{align}
Remarkably, the sum on the right-hand side can be evaluated using a trick, namely averaging over cyclic permutations of $\mbb$.
Note that for $\mbb \in \mcB$,
\begin{equation} \label{eq:exponents-in-f-of-b-sum-to-zero}
    (b_M - 1) + \dots + (b_J - 1) = 0,
\end{equation}
since we have $k$ summands and $b_M + \dots + b_J = k$.
The key tool is the next lemma, which essentially tells us that on average over cyclic permutations, $1/f(\mbb) = 1/k$.

\begin{lemma}[Cycle lemma]
    Suppose $x_1, \dots, x_r$ are positive real numbers with $x_1 \cdots x_r = 1$. 
    Then, reading indices modulo $r$, we have
    \begin{equation*}
        \sum \limits_{t = 0}^{r - 1} \left( \sum \limits_{j = 1}^r x_{t + 1} \cdots x_{t + j} \right)^{-1} = 1.
    \end{equation*}
    More generally, if $x_1 \cdots x_r = X$, then the above sum lies in $[\max(1, X)^{-1}, \min(1, X)^{-1}]$.
\end{lemma}

For the (straightforward) proof, we refer the reader to \cite{integerswithadivisor}.
\begin{remark}
    In \cite{integerswithadivisor}, Ford remarks that the above lemma is related to the following observation.
    Suppose $z_1, \dots, z_r$ are real numbers with $z_1+ \dots + z_r = 0$.
    Then there exists a cyclic permutation of $(z_1, \dots, z_r)$ such that all partial sums are nonnegative, i.e., there exists some $t \in \Set{0, \dots, r - 1}$ such that
    \begin{equation*}
        z_{t+1} + \dots + z_{t + j} \ge 0
    \end{equation*}
    for all $j \ge 1$. This can be seen by choosing $t$ so that $z_1 + \dots + z_t$ is minimal.
    This observation, in turn, is related to a more general \defn{cycle lemma}, which can be used, e.g. to prove Bertrand's ballot theorem, see \cite{aproblemofarrangements}.

    With the above in mind, the idea to average over cyclic permutations can be motivated as follows. 
    We expect 
    \begin{equation*}
        f(\mbb) \approx g(\mbb) = \max_{j} 2^{(b_M - 1) + \dots + (b_j - 1)}.
    \end{equation*}
    Since $(b_M - 1) + \dots + (b_J - 1) = 0$, there exists a cyclic permutation $\mbb'$ of $(b_M, \dots, b_J)$ such that all partial sums are $\le 0$, and so $g(\mbb') = 1$.
    Therefore, we expect that on average over cyclic permutations $\mbb'$ of $\mbb$, $1/f(\mbb') \gg 1/k$.
    Since the product $b_M! \dots b_J!$ is invariant under cyclic permutations, and since the orbits under cyclic permutations partition $\mcB$, we thus expect $1/f(\mbb) \gg 1/k$ on average over $\mbb \in \mcB$.
\end{remark}

Returning to the proof of Theorem \ref{thm:integers-with-a-divisor-in-y-2y}, we have
\begin{equation} \label{eq:estimate-sum-over-b-by-T-and-S}
    \sum \limits_{\mbb \in \mcB^*} \frac{1}{b_M! \cdots b_J! f(\mbb)} \ge S - \sum \limits_{M\le h \le k/M} T(h),
\end{equation}
where 
\begin{align*}
    S = \sum \limits_{\mbb \in \mcB}\frac{1}{b_M! \cdots b_J! f(\mbb)}, \qquad T(h) = \sum \limits_{\substack{\mbb \in \mcB\\b_h > Mh}}\frac{1}{b_M! \cdots b_J!f(\mbb)}.
\end{align*}
Note that for $h > k/M$ the condition $b_h \le Mh$ is automatically satisfied since there are only $k$ prime factors in total.

The cycle lemma can be used to explicitly (!) evaluate $S$ and to show that the total size of the $T(h)$ is acceptable.
\begin{lemma} \label{lem:evaluate-S-and-T}
    We have
    \begin{equation*}
        S = \frac{k^{k - 1}}{k!}
    \end{equation*}
    and for all $M \le h \le k/M$,
    \begin{equation*}
        T(h) \le \frac{2}{(Mh)!} \frac{k^{k - 1}}{k!}.
    \end{equation*}
\end{lemma}

\begin{proof}
    For the first part, we want to apply the cycle lemma with $x_j = 2^{b_{M - 1 + j} - 1}$ ($j \in [k]$).
    Note that \eqref{eq:exponents-in-f-of-b-sum-to-zero} implies $x_1 \cdots x_k = 1$.
    Since every cyclic permutation of $(b_M, \dots, b_J)$ gives a 1-to-1 map on $\mcB$, and since the product of factorials is invariant under cyclic permutations, we can write
    \begin{equation*}
        \sum \limits_{\mbb \in \mcB} \frac{1}{b_M! \cdots b_J! f(\mbb)} = \frac{1}{k} \sum \limits_{\mbb \in \mcB} \frac{1}{b_M! \cdots b_J!} \sum \limits_{\mbb'} \frac{1}{f(\mbb')},
    \end{equation*}
    where $\mbb'$ runs over all $k$ cyclic permutations of $(b_M, \dots, b_J)$.
    By the cycle lemma, the sum over $\mbb'$ equals $1$. 
    Using the multinomial theorem, we conclude that
    \begin{align*}
        S &= \frac{1}{k} \sum \limits_{\mbb \in \mcB} \frac{1}{b_M! \cdots b_J!} = \frac{1}{k \cdot k!} \sum \limits_{\mbb \in \mcB} \binom{k}{b_M, \dots, b_J} = \frac{k^{k - 1}}{k!}.
    \end{align*}

    The second part follows a similar pattern but is technically more involved. The idea is to apply the cycle lemma to cyclic permutations of $(b_{h + 1}, \dots, b_J)$.
    Details can be found in \cite{fordlecturenotes}.
\end{proof}

We finally have collected everything we need to complete the proof of the lower bound.
\begin{proof}[Proof of Theorem \ref{thm:integers-with-a-divisor-in-y-2y}, lower bound]
    If $y \le y_0$, then $H(x, y, 2y) \asymp x$.
    Hence, we may assume that $y$ is large.
    Lemma \ref{lem:evaluate-S-and-T} gives
    \begin{equation*}
        \sum \limits_{M \le h \le k/M} T(h) \le \sum \limits_{h \ge M}\frac{2}{(Mh)!} \frac{k^{k - 1}}{k!} \le \frac{1}{2} \frac{k^{k - 1}}{k!},
    \end{equation*}
    and so 
    \begin{equation*}
        S - \sum \limits_{M \le h \le k/M} T(h) \ge \frac{1}{2} \frac{k^{k - 1}}{k!},
    \end{equation*}
    provided that $M$ is large enough.
    Combining \eqref{eq:estimate-h-x-y-2y-by-sum-over-f-b} and \eqref{eq:estimate-sum-over-b-by-T-and-S}, we deduce that
    \begin{align} \label{eq:final-calculation-lower-bound}
        H(x, y, 2y) &\gg \frac{x(2 \log 2)^k}{(\log y)^2} \sum \limits_{\mbb \in \mcB^*} \frac{1}{b_M! \cdots b_J! f(\mbb)} \nonumber \\
            &\gg \frac{x(2 \log 2)^k}{(\log y)^2} \left( S - \sum \limits_{M \le h \le k/M} T(h) \right) \nonumber \\ 
            &\gg \frac{x(2 \log 2)^k}{(\log y)^2} \frac{k^{k - 1}}{k!}.
    \end{align}
    Stirling's formula states that $
        k! \sim \sqrt{2 \pi k} (k/e)^k$.
    Substituting this into \eqref{eq:final-calculation-lower-bound} and using that $k = \lflr \log \log y / \log 2 \rflr + O(1)$, we obtain
    \begin{align*}
        H(x, y, 2y) &\gg \frac{x}{(\log y)^2 (\log \log y)^{3/2}} \exp \left( \frac{1 + \log 2 + \log \log 2}{\log 2} \cdot \log \log y \right) \\ 
        &\gg \frac{x}{(\log y)^\delta (\log \log y)^{3/2}},
    \end{align*}
    as desired.
\end{proof}

\subsection{Proof of the upper bound}
In this section, we prove the upper bound for
\begin{equation} \label{eq:sum-to-be-upper-bounded}
    H(x, y, 2y) \ll \frac{x}{(\log y)^2} \sum \limits_{a \in \mcS(y)} \frac{L(a)}{a}.
\end{equation}
To this end, we split the sum according to $\omega(a)$. 
Lemma \ref{lem:properties-of-L} (i) tells us that when $\omega(a) = k$, then 
\begin{equation*}
    L(a) \le \min\Set{\tau(a) \log 2,\, \log (2a)} \lessapprox \min \Set{2^k,\, \log y}.
\end{equation*}
So again, we expect a transition at the point where $2^k = \log y$, i.e., at 
\begin{equation*}
    k_0 = \lflr \frac{\log \log y}{\log 2} \rflr.
\end{equation*}
Define
\begin{equation*}
    S_k(y) = \sum \limits_{\substack{a \in \mcS(y)\\\omega(a) = k}} \frac{L(a)}{a}.
\end{equation*}
In light of \eqref{eq:sum-to-be-upper-bounded}, we want to show that 
\begin{equation*}
    \sum \limits_{k} S_k(y) \ll \frac{(\log y)^{2- \delta}}{(\log \log y)^{3/2}}
\end{equation*}
in order to complete the proof of Theorem \ref{thm:integers-with-a-divisor-in-y-2y}.
When $|k - k_0|$ is relatively large, the simple bounds from Lemma \ref{lem:properties-of-L} (i) suffice to show that the contribution of $S_k(y)$ to \eqref{eq:sum-to-be-upper-bounded} is negligible. 
The main contribution comes from integers with approximately $k_0$ prime factors.
In this case, a more sophisticated analysis is necessary, bounding $S_k(y)$ in terms of a multivariate integral and using an estimate about uniform order statistics. 

We start with the case where $|k - k_0|$ is large.
\begin{lemma} \label{lem:bound-case-where-k-k0-is-large}
    We have 
    \begin{equation*}
        \sum \limits_{k\colon \left|k - k_0\right| \ge k_0/2} S_k(y) \ll (\log y)^{3/2}.
    \end{equation*}
\end{lemma}

This is much less than the alleged main term $(\log y)^{2 - \delta} / (\log \log y)^{3/2}$ (since \mbox{$\delta < 0.1$}).

\begin{proof}
    When $k \le k_0/2$, we can use the bound $L(a) \le 2^k$. 
    Then 
    \begin{equation*}
        S_k(y) \le 2^k \sum \limits_{\substack{a \in \mcS(y) \\\omega(a) = k}} \frac{1}{a} \le \frac{2^k}{k!} \left( \sum \limits_{p \le y} \frac{1}{p} \right)^k \ll \frac{(2 \log \log y)^k}{k!}.
    \end{equation*}
    We have $2 \log \log y /k \ge 4 \log 2 > 2$, so
    \begin{align*}
        \sum \limits_{k \le k_0/2} S_k(y) &\ll \sum \limits_{k \le k_0/2} 2^{k - k_0} \frac{{(2 \log \log y)}^{k_0}}{k_0!} \\ 
        &\ll 2^{-k_0/2} \frac{(\log y)^{2 - \delta}}{\sqrt{\log \log y}} \\
        &\ll (\log y)^{3/2}.
    \end{align*}
    Here, the second inequality follows from Stirling's formula. 
    For the last inequality, we used that $2^{-k_0/2} \le (\log y)^{-1/2}$ and $\delta > 0$.

    When $k \ge \frac{3}{2}k_0$, we use the bound $L(a) \le \log 2a$.
    We have
    \begin{align*}
        S_k(y) &\ll \sum\limits_{\substack{a \in \mcS(y) \\\omega(a) = k}} \frac{1}{a} \sum \limits_{p \mid a} \log p \\ 
            &\le \left(\sum \limits_{p \le y} \frac{\log p}{p}\right) \left(\sum \limits_{\substack{b \in \mcS(y)\\ \omega(b) = k - 1}} \frac{1}{b} \right) \\ 
            &\ll \log y \cdot \frac{(\log \log y + O(1))^{k - 1}}{(k - 1)!},
    \end{align*}
    where we used Mertens' first and second estimates in the last step.
    Hence, 
    \begin{equation*}
        \sum \limits_{k \ge 3k_0 / 2} S_k(y) \ll (\log y) \sum \limits_{k \ge 3k_0/ 2} \frac{(\log \log y + O(1))^{k - 1}}{(k - 1)!}.
    \end{equation*}
    The sum decays at least geometrically.
    Bounding it by its largest term and applying Stirling's formula as above, we obtain 
    \begin{equation*}
        \sum \limits_{k \ge 3/2 k_0} S_k(y) \ll (\log y)^{3/2},
    \end{equation*}
    as claimed.
\end{proof}

It remains to bound the contribution of integers with roughly $k_0$ prime divisors. 
We could try the same argument as above, but this would only give a bound of $(\log y)^{2 - \delta} / (\log \log y)^{1/2}$, which is too large by a factor of $\log \log y$.
The reason is that the bound $L(a) \le \min(\tau(a) \log 2, \log(2a))$ is only close to being sharp when the divisors of $a$ are evenly distributed on a logarithmic scale.

To obtain the correct bound, one must bound $L(a)$ more carefully.
One reason for the divisors of $a$ to cluster is when one of the prime factors of $a$ is significantly smaller than expected. 
This is captured by the third bound in Lemma \ref{lem:properties-of-L}, which states that for $p_1 < \dots < p_k$,
\begin{equation} \label{eq:recall-properties-of-L-3}
    L(p_1 \cdots p_k) \le \min 2^{k-j} \cdot \log (2p_1 \cdots p_j).
\end{equation}
Define
\begin{equation*}
    F(\mbt) = \min \limits_{1 \le j \le k} 2^{-j} \left( 1 + 2^{t_1} + \dots + 2^{t_j} \right).
\end{equation*}
With this notation, \eqref{eq:recall-properties-of-L-3} translates to
\begin{equation*}
    L(p_1 \cdots p_k) \le 2^{k}\cdot F\left( \frac{\log \log p_1}{\log 2}, \dots, \frac{\log \log p_k}{\log 2} \right).
\end{equation*}
The definition of $F$ is motivated by the idea that looking at prime divisors on a log-log scale leads to the connection to uniform order statistics mentioned earlier.

As a general principle, if $f: \mR \rightarrow \mR$ is some suitably nice function,
we expect 
\begin{equation*}
    \sum \limits_{p \le y} \frac{f(\log \log p)}{p} \approx \int_{2}^y \frac{f(\log \log t)}{t} \frac{\mathrm d t}{\log t},
\end{equation*}
where the factor $1/\log t$ accounts for the density of the primes around $t$.
Changing variables to $\xi = \log \log t / \log \log y$, this is
\begin{equation*}
    \approx \log \log y \int_0^1 f(\xi) \,\mathrm d \xi.
\end{equation*}
Applying this intuition $k$ times to $F$, we would expect that
\begin{align*}
    \sum \limits_{p_1 < p_2 < \dots < p_k \le y} & \frac{F\left( \frac{\log \log p_1}{\log 2}, \dots, \frac{\log \log p_k}{\log 2} \right)}{p_1 \cdots p_k} \\
        & \approx\; (\log \log y)^k  \idotsint \limits_{0 \le \xi_1 \le \dots \le \xi_k\le 1} F(\xi_1 k_0, \dots, \xi_k k_0) \; \mathrm{d} \bm{\xi}.
\end{align*}
Unfortunately, $F$ is not a particularly nice function, so showing something like the above is not a straightforward matter of applying partial summation. 
However, one can exploit the monotonicity of $F$ to give an upper bound of the same shape, which is what the next lemma achieves.

\begin{lemma} \label{lem:approximate-sum-by-integral}
    Suppose $\frac{k_0}{2} \le k \le 2k_0$. 
    Then 
    \begin{equation*}
        S_k(y) \ll (2 \log \log y)^k \cdot U_k(k_0),
    \end{equation*}
    where 
    \begin{equation*}
        U_k(s) = \idotsint \limits_{0 \le \xi_1 \le \dots \le \xi_k\le 1} \min \limits_{0 \le j \le k} 2^{-j} \left( 1 + 2^{s\xi_1} + \dots + 2^{s \xi_j} \right) \; \mathrm{d} \bm{\xi}.
    \end{equation*}
\end{lemma}

\begin{remark}
    The quantity $U_k(s)$ has a natural interpretation in terms of uniform order statistics. 
    Suppose $X_1, \dots, X_k$ are i.i.d. random variables, uniformly distributed on $[0, 1]$, and let $0 \le \xi_1 \le \dots \le \xi_k \le 1$ denote their order statistics.
    Then 
    \begin{equation*}
        U_k(s) = \frac{1}{k!}\expn{\min \limits_{0 \le j \le k} 2^{-j}\left( 1 + 2^{s\xi_1} + \dots + 2^{s \xi_j} \right)}.
    \end{equation*}
    Typically, $\xi_j \approx j/(k + 1)$. 
    When $s$ is close to $k$, we expect
    \begin{equation*}
        \min_{0 \le j \le k} 2^{-j}\left( 1 + 2^{s\xi_1} + \dots + 2^{s \xi_j} \right) \approx \min_{1 \le j \le k} 2^{s\xi_j - j}.
    \end{equation*}
    This is small when one of the $\xi_j$ is significantly smaller than expected, so $U_k(s)$ essentially measures how well-spread the $\xi_1, \dots, \xi_k$ are. 
    This connects back to the discussion in Section \ref{subsec:heuristic} and \eqref{eq:heuristic-condition-on-order-statistics} specifically.
\end{remark}

Before turning to the proof of Lemma \ref{lem:approximate-sum-by-integral}, we record that
\begin{equation} \label{eq:volume-of-ordering}
    \Vol \set{\mathbf{x} \in \mR^k}{0 \le x_1 \le \dots \le x_k \le 1} = \frac{1}{k!}.
\end{equation}
This can be seen by noting that the regions corresponding to the $k!$ different orderings of $x_1, \dots, x_k$ partition the hypercube $[0, 1]^k$, and by symmetry, all these regions must have the same volume.
%Alternatively, one can also compute the volume by just integrating directly.

Similarly to the proof of Lemma \ref{lem:upper-bound-for-sum-over-W-a}, the idea is to localise prime factors.
Recall the definition of $\lambda_i$ and $D_i$ from Section \ref{subsec:proof-of-lower-bound}: $D_i = (\lambda_{i - 1}, \lambda_i]$, where the $\lambda_i$ are defined inductively by $\lambda_0 = 3/2$ and, for $i \ge 1$, $\lambda_i$ is the largest prime such that 
\begin{equation*}
    \sum \limits_{\lambda_{i - 1} < p \le \lambda_i} \frac{1}{p} \le \log 2.
\end{equation*}

\begin{proof}[Proof of Lemma \ref{lem:approximate-sum-by-integral}]
    Consider integers $a = p_1 \cdots p_k$ with $p_1 < \dots < p_k \le y$.
    For $\ell \in [k]$, let $j_\ell$ be such that $p_\ell \in D_{j_\ell}$.
    Recall that $p \in D_{j}$ implies that $\log p \le \log \lambda_{j} \le 2^{j + O(1)}$.
    Therefore, Lemma \ref{lem:properties-of-L} (iii) gives
    \begin{align*}
        L(a) &\le 2^k \min \limits_{0 \le \ell \le k} 2^{-\ell} \log (2 p_1 \cdots p_\ell) \\
            &\ll 2^k \min \limits_{0 \le \ell \le k} 2^{-\ell} \left(1 + 2^{j_1} + \dots + 2^{j_\ell} \right) \\
            &= 2^k \cdot F(j_1, \dots, j_k).
    \end{align*}
    Next, we have $\log \log p_\ell = (\log 2) j_\ell + O(1)$. 
    Since also $p_\ell \le y$, we can bound $j_\ell \le K$, where
    \begin{equation*}
        K = k_0 + O(1) = \lflr \frac{\log \log y}{\log 2} \rflr + O(1).
    \end{equation*}
    Hence, 
    \begin{equation*}
        S_k(y) \;=\; \sum \limits_{\substack{a \in \mcS(y) \\ \omega(a) = k}} \frac{L(a)}{a} 
            \;\ll\; 2^k \sum \limits_{1 \le j_1 \le \dots \le j_k \le K} F(\mbj) \sum \limits_{\substack{p_1 < \dots < p_k \\ p_\ell \in D_{j_\ell} \; \forall \ell \in [k]}} \frac{1}{p_1 \cdots p_k}.
    \end{equation*}
    Now consider the contribution of the innermost sum for a fixed vector $\mbj = (j_1, \dots, j_k)$ with $1 \le j_1 \le \dots \le j_k \le K$.
    Write $b_m$ for the number of primes in $D_m$ ($1 \le m \le K$), that is, $b_m$ is the number of entries $j_\ell$ of $\mbj$ with $j_\ell = m$.
    Then $b_1 + \dots + b_K = k$, and we have
    \begin{equation*}
        \sum \limits_{\substack{p_1 < \dots < p_k \\ p_\ell \in D_{j_\ell} \; \forall \ell \in [k]}} \frac{1}{p_1 \cdots p_k} 
            \;\le\; \prod_{m = 1}^K \frac{1}{b_m!} \left(  \sum \limits_{p \in D_m} \frac{1}{p} \right)^{b_m} 
            \;\le \; \prod\limits_{m = 1}^K \frac{(\log 2)^{b_m}}{b_m!} 
            \; = \; \frac{(\log 2)^k}{b_1! \cdots b_K!}.
    \end{equation*}
    By \eqref{eq:volume-of-ordering}, we can write 
    \begin{equation*}
        \frac{(\log 2)^k}{b_1! \cdots b_K!} \;=\; (\log 2)^k \prod\limits_{m = 1}^K\;\; \idotsint \limits_{m - 1 \le t_1 \le \dots \le t_{b_m} \le m} 1\, \mathrm{d} \mbt \;=\; (\log 2)^k \int_{R(\mbj)} 1\, \mathrm d \mbt,
    \end{equation*}
    where 
    \begin{equation*}
        R(\mbj) = \set{\mbt \in \mR^{k}}{m - 1 \le t_{B_{m - 1} + 1} \le \dots \le t_{B_m} \le m \quad \forall m \in [K]}
    \end{equation*}
    and $B_m = b_1 + \dots + b_m$.
    Putting everything together, we arrive at
    \begin{equation*}
        S_k(y) \ll (2 \log 2)^k \sum \limits_{1 \le j_1 \le \dots \le j_k \le K} F(\mbj) \int_{R(\mbj)} 1\, \mathrm d \mbt.
    \end{equation*}
    We now allow $F$ to take values in $\mR^k_{\ge 0}$.
    Clearly, if $t_\ell \in [j_\ell - 1, j_\ell]$ for all $1 \le \ell \le k$, then $F(j_1, \dots, j_k) \le 2 F(t_1, \dots, t_k)$.
    Thus, 
    \begin{align*}
        S_k(y) &\ll (2 \log 2)^k \sum \limits_{1 \le j_1 \le \dots \le j_k \le K}  \int_{R(\mbj)} F(\mbt)\, \mathrm d \mbt\\
            &\le (2 \log 2)^k \idotsint\limits_{0 \le t_1 \le \dots \le t_k \le K} F(\mbt)\, \mathrm{d} \mbt.
    \end{align*}
    Changing variables to $\xi_i = t_i / K$, we conclude that
    \begin{align*}
        S_k(y) &\ll (2 \log 2)^k K^k \idotsint\limits_{0 \le \xi_1 \le \dots \le \xi_k \le 1} F(K \xi_1, \dots, K\xi_k) \, \mathrm{d} \bm{\xi} \\
            &\ll ((2 \log 2) k_0 + O(1))^k  \cdot U_k(k_0) %\idotsint\limits_{0 \le \xi_1 \le \dots \le \xi_k \le 1} F(K \xi_1, \dots, K\xi_k) \, \mathrm{d} \bm{\xi}
    \end{align*}
    Finally, $((2 \log 2) k_0 + O(1))^k \ll (2 \log \log y)^k$ since $k \le 2k_0$. 
    This concludes the proof.
\end{proof}

To conclude the proof of the upper bound, we now need to estimate 
\begin{equation*}
    U_k(s) =\idotsint \limits_{0 \le \xi_1 \le \dots \le \xi_k\le 1} \min \limits_{0 \le j \le k} 2^{-j} \left( 1 + 2^{s\xi_1} + \dots + 2^{s \xi_j} \right) \; \mathrm{d} \bm{\xi}.
\end{equation*} 
The region of integration in $U_k(s)$ has volume $1/k!$, and the integrand is always $\ll \min (1, k 2^{s-k})$ (since we can always take $j = 0$ or $j=k$ in the definition of $U_k(s)$).
Therefore, we have the trivial bound
\begin{equation*}
    U_k(s) \ll \frac{\min (1, k\cdot2^{s - k})}{k!} \ll \frac{1}{(k - 1)! (2^{k - s} + 1)}. 
\end{equation*}
Typically (e.g., when $\xi_\ell \approx \ell/(k + 1)$), the integrand should only be roughly $\ll \min(1, 2^{s - k})$,
so we would perhaps expect to obtain a bound like
\begin{equation*}
    U_k(s) \lessapprox \frac{1}{k!(2^{k - s} + 1)}.
\end{equation*}
The following lemma improves upon this by another factor of $1/(k + 1)$, at least when $|s - k|$ is not too large.
This will be crucial in obtaining the correct exponent for the $\log \log y$ term in the upper bound for $H(x, y, 2y)$.
\begin{lemma}[{\cite[Lemma 3.6]{integerswithadivisor}}]
    \label{lem:estimate-U-k-s}
    Uniformly for $0 \le k \le 10 s$, we have
    \begin{equation*}
        U_k(s) \ll \frac{1 + |s - k|^2}{(k + 1)!( 2^{k - s} + 1)}.
    \end{equation*}
\end{lemma}

The proof of Lemma \ref{lem:estimate-U-k-s} is technically involved, and we refer the reader to \cite{integerswithadivisor} for details.
However, assuming the lemma, the proof of the upper bound is straightforward. 
\begin{proof}[Proof of Theorem \ref{thm:integers-with-a-divisor-in-y-2y}, upper bound]
    We start by bounding the contribution of $k \ge k_0$. Observe that in this case,
    \begin{equation*}
        k \ge \lflr\frac{\log \log y}{\log 2} \rflr \ge \frac{4}{3} \log \log y.
    \end{equation*}
    Therefore, combining Lemmas \ref{lem:approximate-sum-by-integral} and \ref{lem:estimate-U-k-s}, we have
    \begin{align*}
        \sum \limits_{k_0 \le k \le 2k_0} S_k(y) &\ll \sum \limits_{k_0 \le k \le 2k_0} (2 \log \log y)^k U_k (k_0) \\
        &\ll \sum \limits_{k_0 \le k \le 2k_0} (2\log \log y)^k \frac{1 + (k - k_0)^2}{(k + 1)! (2^{k - k_0} + 1)} \\ 
        &\ll 2^{k_0} \frac{(\log \log y)^{k_0}}{(k_0 + 1)!} \sum \limits_{k_0 \le k \le 2k_0} \left(1 + (k - k_0)^2 \right) \frac{(k_0 + 1)!(\log \log y)^{k-k_0}}{(k + 1)!} \\
        & \ll 2^{k_0} \frac{(\log \log y)^{k_0}}{(k_0 + 1)!} \sum \limits_{m \ge 0} \frac{1 + m^2}{\left( \frac{4}{3} \right)^m}.
    \end{align*}
    The sum over $m$ in the last line is clearly $O(1)$.
    We conclude that
    \begin{align*}
        \sum \limits_{k_0 \le k \le 2k_0} S_k(y) &\ll \frac{(2 \log \log y)^{k_0}}{(k_0 + 1)!} \ll \frac{(\log y)^{2 - \delta}}{(\log \log y)^{3/2}},
    \end{align*}
    where the last inequality follows from Stirling's formula.

    A similar argument gives that
    \begin{equation*}
        \sum \limits_{k_0/2 \le k \le k_0} S_k(y) \ll\frac{(\log y)^{2 - \delta}}{(\log \log y)^{3/2}}.
    \end{equation*}

    We have already proven in Lemma \ref{lem:bound-case-where-k-k0-is-large} that the contribution of $\left|k - k_0\right| \ge k_0 / 2$ is $\ll (\log y)^{3/2}$, which is negligible compared to the main term. 
    Therefore, the proof of the upper bound is complete.
\end{proof}
